% !PRINTED 40
\weberpage{40}
For if $\alpha$ and $\beta$ are given arbitrarily, one can determine $s$ with respect to the modulus $ab$ so that
\[
s \equiv \alpha \ (\mathrm{mod}\ a), \quad s \equiv \beta \ (\mathrm{mod}\ b),
\]
whereby (2) becomes identical with (1). Nevertheless, in a certain sense the constitution of the basis is completely determined by the nature of the groups, according to the following theorem:

\gesperrt{II. If
\[
A_1, A_2 \ldots A_\nu
\]
with the degrees
\[
a_1, a_2 \ldots a_\nu
\]
and
\[
B_1, B_2 \ldots B_\mu
\]
with the degrees
\[
b_1, b_2 \ldots b_\mu
\]
are two bases of an Abelian group $S$ of degree $n$, if $p$ is a prime dividing $n$ and
\weq{1}{a_1 = p_1 a_1', \quad a_2 = p_2 a_2', \quad \ldots \quad a_\nu = p_\nu a_\nu',}
\weq{2}{b_1 = p_1' b_1', \quad b_2 = p_2' b_2', \quad \ldots \quad b_\mu = p_\mu' b_\mu',}
where $p_1, p_2 \ldots p_\nu, p_1', p_2' \ldots p_\mu'$ are the highest powers of $p$ that divide $a_1, a_2 \ldots a_\nu, b_1, b_2 \ldots b_\mu$, then all the prime powers $p_1, p_2 \ldots p_\nu$ that are greater than $1$ also occur among the $p_1', p_2' \ldots p_\mu'$, and conversely.}

To prove this theorem, we take the elements of the two bases $A$ and $B$ to be so ordered that
\weq{3}{\begin{aligned}
p_1 &\geqq p_2 \geqq \ldots \geqq p_\nu \\
p_1' &\geqq p_2' \geqq \ldots \geqq p_\mu'.
\end{aligned}}
The integers $a_1', a_2' \ldots a_\nu', b_1', b_2' \ldots b_\mu'$ occurring in (1) and (2) are by their definition not divisible by $p$, and if we therefore denote by $m$ the least common multiple of $a_1', a_2' \ldots a_\nu'$, then $m$ too is not divisible by $p$. But if
\weq{4}{\Theta = A_1^{\alpha_1} A_2^{\alpha_2} \ldots A_\nu^{\alpha_\nu}}
is an arbitrary element of $S$, it follows that
\[
\Theta^{p_1 m} = 1
\]
must hold.
